% =====================================================================
%  APÊNDICE B — CONFIGURAÇÃO DO AMBIENTE
% =====================================================================

\chapter{CONFIGURAÇÃO DO AMBIENTE}
\label{ap:ambiente}

Este apêndice descreve a configuração dos principais componentes do ambiente de
execução do \argus{}.

\section{Conexão com sensores Bluetooth Low Energy}
\label{ap:ambiente-ble}

Os sensores BLE (cinta cardíaca, medidor de potência e rolo inteligente) são
detectados automaticamente pelo escaneamento do sistema. No Windows, caso a
pilha Bluetooth não exponha os UUIDs de serviço durante o escaneamento, o
\argus{} lista os dispositivos nomeados e valida os serviços FTMS/HR após a
seleção e conexão. Para diagnóstico, a variável de ambiente
\code{ARGUS\_BLE\_DEBUG=1} ativa logs detalhados de escaneamento.

\complementar{Descrever aqui, se necessário, a configuração específica do
\code{virtual trainer} e o procedimento de pareamento dos sensores nos
sistemas operacionais suportados (Windows, Linux e Android).}

\section{Assistente de treino com IA local (Ollama)}
\label{ap:ambiente-ollama}

O assistente de IA requer uma instância do Ollama \cite{ollama2024} executando
no dispositivo. Recomenda-se o modelo quantizado de pequena escala
\texttt{qwen2.5:3b}, que executa de forma responsiva em CPU. A seleção do
modelo é feita na interface do \argus{} (\textit{AIChatController}).

\complementar{Descrever aqui o passo a passo de instalação do Ollama
(\url{https://ollama.com/}), o comando de \textit{pull} do modelo
recomendado e a configuração da variável de ambiente do servidor Ollama, caso
haja necessidade de apontar para um \textit{endpoint} diferente do padrão.}

\section{Integração com o Strava e o Intervals.icu}
\label{ap:ambiente-strava}

As integrações com o Strava e o Intervals.icu exigem o cadastro de uma
aplicação no respectivo portal de desenvolvedores e o armazenamento das
credenciais de OAuth 2.0. O fluxo de autorização é realizado por um servidor
HTTP local, e o \textit{refresh token} é gerenciado automaticamente pelo
sistema.

\complementar{Descrever aqui as credenciais necessárias (client ID e client
secret) e o procedimento de obtenção junto à Strava
(\url{https://developers.strava.com/}) e ao Intervals.icu
(\url{https://intervals.icu/}), bem como as permissões solicitadas ao usuário.}
